A GPS-denied UAV can navigate by matching the terrain heights measured by its barometric and
radar altimeters against an elevation map. This report studies that problem --- terrain-aided
navigation --- in simulation over real Copernicus GLO-30 terrain of three contrasting regions of
Sri Lanka, chosen by a pre-registered rule. A vectorised sequential importance resampling
particle filter is compared with an extended Kalman filter and INS dead reckoning in seven
Monte Carlo experiments with common random numbers, paired tests and confidence intervals. On real
terrain the EKF's linearisation fails earlier than on smooth synthetic hills: over the highlands it
converges in only \EOneHillSFiveZeroZeroEkfConv\,\% of runs at 500\,m initial uncertainty (synthetic
terrain: \EOneSynthSFiveZeroZeroEkfConv\,\%) and, without innovation gating, ends worse than dead
reckoning, while a 20\,000-particle filter converges in
\EOneHillSFiveZeroZeroPfTwoZeroZeroZeroZeroConv\,\% and tracks at about
\EOneHillSFiveZeroZeroPfTwoZeroZeroZeroZeroTrack\,m. Across terrain types, rougher terrain lowers the
particle filter's error and raises the EKF's. The particle filter's characteristic failure, false
lock onto a wrong patch of terrain, is countered by roughening and by an online divergence monitor
based on the normalised innovation squared: together they raise convergence at 2\,km initial
uncertainty from \ESixHillSTwoZeroZeroZeroPfFiveZeroZeroZeroConv\,\% to
\ESixHillSTwoZeroZeroZeroPfMonitorRoughConv\,\%. A Student-$t$ likelihood keeps the filter robust to
radar outliers, and a Rao-Blackwellised filter estimates the INS bias much better once the position
is locked on. In the tracking phase the 20\,000-particle filter stays about twice above the
posterior Cramér--Rao bound, a gap that a mismatched-gain analysis attributes almost entirely to its
deliberately inflated process noise.
